%% ---------------------------------------------------------------------
%% 11. ASSESSMENT AND ACCEPTANCE CRITERIA — Required · ½–1 page
%% How the work will be judged — so the student knows what "good" means.
%% Use the rubric table: 4–7 criteria, weights summing to 100%, and for each
%% criterion a concrete description of minimum (pass), good and excellent.
%% Also list hard acceptance criteria (must all be met for a pass).
%% Align with the university's official grading rules; this does not replace them.
%% ---------------------------------------------------------------------
\chapter{Assessment and Acceptance Criteria}
\label{ch:assessment}

\section{Acceptance criteria}
The project is accepted for defense when:
\begin{itemize}
  \item the MVP (Chapter~\ref{ch:scope}) works and is demonstrated;
  \item \missing{criterion}
\end{itemize}

\section{Assessment rubric}
\begin{rubric}[Assessment rubric.]
  \criterion{Technical result}{35\%}{MVP works}{MVP + measured performance}{\missing{MVP + extension, clean design}}
  \criterion{Verification and testing}{20\%}{\missing{…}}{\missing{…}}{\missing{…}}
  \criterion{Understanding and analysis}{15\%}{\missing{…}}{\missing{…}}{\missing{…}}
  \criterion{Independence and work process}{10\%}{\missing{…}}{\missing{…}}{\missing{…}}
  \criterion{Document and presentation}{20\%}{\missing{…}}{\missing{…}}{\missing{…}}
\end{rubric}
